%% Copyright (C) 2020 Saurabh Joshi
%% 
\let\negmedspace\undefined
\let\negthickspace\undefined

\documentclass[journal,12pt,onecolumn]{IEEEtran}
%\documentclass[journal,12pt,twocolumn]{IEEEtran}
%
\usepackage{setspace}
\usepackage{gensymb}
%\doublespacing
\singlespacing

%\usepackage{graphicx}
%\usepackage{amssymb}
%\usepackage{relsize}
\usepackage[cmex10]{amsmath}
%\usepackage{amsthm}
%\interdisplaylinepenalty=2500
%\savesymbol{iint}
%\usepackage{txfonts}
%\restoresymbol{TXF}{iint}
%\usepackage{wasysym}
\usepackage{amsthm}
\usepackage{mathrsfs}
\usepackage{txfonts}
\usepackage{stfloats}
\usepackage{cite}
\usepackage{cases}
\usepackage{subfig}
%\usepackage{xtab}
\usepackage{longtable}
\usepackage{multirow}
%\usepackage{algorithm}
%\usepackage{algpseudocode}
\usepackage{enumitem}
\usepackage{mathtools}
\usepackage{tikz}
\usetikzlibrary{shapes,arrows,positioning}
\usepackage{circuitikz}
\usepackage{verbatim}
\usepackage{hyperref}
%\usepackage{stmaryrd}
\usepackage{tkz-euclide} % loads  TikZ and tkz-base
%\usetkzobj{all}
\usepackage{listings}
    \usepackage{color}                                            %%
    \usepackage{array}                                            %%
    \usepackage{longtable}                                        %%
    \usepackage{calc}                                             %%
    \usepackage{multirow}                                         %%
    \usepackage{hhline}                                           %%
    \usepackage{ifthen}                                           %%
  %optionally (for landscape tables embedded in another document): %%
    \usepackage{lscape}     
\usepackage{multicol}
\usepackage{chngcntr}
\usepackage{iftex}
%\usepackage[latin9]{inputenc}
\usepackage{geometry}
%\geometry{verbose,tmargin=2cm,bmargin=3cm,lmargin=1.8cm,rmargin=1.5cm,headheight=2cm,headsep=2cm,footskip=3cm}
\usepackage{array}
\newcolumntype{L}[1]{>{\raggedright\let\newline\\\arraybackslash\hspace{0pt}}m{#1}}
\newcolumntype{C}[1]{>{\centering\let\newline\\\arraybackslash\hspace{0pt}}m{#1}}
\newcolumntype{R}[1]{>{\raggedleft\let\newline\\\arraybackslash\hspace{0pt}}m{#1}}
 \usepackage{float}
%\usepackage{graphicx}
%\usepackage{setspace}
%\usepackage{parskip}

\def \hsp {\hspace{3mm}}

\makeatletter

\providecommand{\tabularnewline}{\\}


\makeatother
\ifxetex
\usepackage[T1]{fontenc}
\usepackage{fontspec}
%\setmainfont[ Path = fonts/]{Sanskrit_2003.ttf}
\newfontfamily\nakulafont[Script=Devanagari,AutoFakeBold=2,Path = fonts/]{Nakula}
%\newfontfamily\liberationfont{Liberation Sans Narrow}
%\newfontfamily\liberationsansfont{Liberation Sans}
\fi
\usepackage{tikz}
\usepackage{xcolor}
%\usepackage{enumerate}

%\usepackage{wasysym}
%\newcounter{MYtempeqncnt}
\DeclareMathOperator*{\Res}{Res}
%\renewcommand{\baselinestretch}{2}
\renewcommand\thesection{\arabic{section}}
\renewcommand\thesubsection{\thesection.\arabic{subsection}}
\renewcommand\thesubsubsection{\thesubsection.\arabic{subsubsection}}

\renewcommand\thesectiondis{\arabic{section}}
\renewcommand\thesubsectiondis{\thesectiondis.\arabic{subsection}}
\renewcommand\thesubsubsectiondis{\thesubsectiondis.\arabic{subsubsection}}

% correct bad hyphenation here
\hyphenation{op-tical net-works semi-conduc-tor}
\def\inputGnumericTable{}                                 %%

\lstset{
language=tex,
frame=single, 
breaklines=true
}

%\begin{document}
%


\newtheorem{theorem}{Theorem}[section]
\newtheorem{problem}{Problem}
\newtheorem{proposition}{Proposition}[section]
\newtheorem{lemma}{Lemma}[section]
\newtheorem{corollary}[theorem]{Corollary}
\newtheorem{example}{Example}[section]
\newtheorem{definition}[problem]{Definition}
%\newtheorem{thm}{Theorem}[section] 
%\newtheorem{defn}[thm]{Definition}
%\newtheorem{algorithm}{Algorithm}[section]
%\newtheorem{cor}{Corollary}
\newcommand{\BEQA}{\begin{eqnarray}}
\newcommand{\EEQA}{\end{eqnarray}}
\newcommand{\define}{\stackrel{\triangle}{=}}

\bibliographystyle{IEEEtran}
%\bibliographystyle{ieeetr}


\providecommand{\mbf}{\mathbf}
\providecommand{\pr}[1]{\ensuremath{\Pr\left(#1\right)}}
\providecommand{\qfunc}[1]{\ensuremath{Q\left(#1\right)}}
\providecommand{\sbrak}[1]{\ensuremath{{}\left[#1\right]}}
\providecommand{\lsbrak}[1]{\ensuremath{{}\left[#1\right.}}
\providecommand{\rsbrak}[1]{\ensuremath{{}\left.#1\right]}}
\providecommand{\brak}[1]{\ensuremath{\left(#1\right)}}
\providecommand{\lbrak}[1]{\ensuremath{\left(#1\right.}}
\providecommand{\rbrak}[1]{\ensuremath{\left.#1\right)}}
\providecommand{\cbrak}[1]{\ensuremath{\left\{#1\right\}}}
\providecommand{\lcbrak}[1]{\ensuremath{\left\{#1\right.}}
\providecommand{\rcbrak}[1]{\ensuremath{\left.#1\right\}}}
\theoremstyle{remark}
\newtheorem{rem}{Remark}
\newcommand{\sgn}{\mathop{\mathrm{sgn}}}
\providecommand{\abs}[1]{\left\vert#1\right\vert}
\providecommand{\res}[1]{\Res\displaylimits_{#1}} 
\providecommand{\norm}[1]{\left\lVert#1\right\rVert}
%\providecommand{\norm}[1]{\lVert#1\rVert}
\providecommand{\mtx}[1]{\mathbf{#1}}
\providecommand{\mean}[1]{E\left[ #1 \right]}
\providecommand{\fourier}{\overset{\mathcal{F}}{ \rightleftharpoons}}
%\providecommand{\hilbert}{\overset{\mathcal{H}}{ \rightleftharpoons}}
%\providecommand{\system}{\overset{\mathcal{H}}{ \longleftrightarrow}}
\providecommand{\system}[1]{\overset{\mathcal{#1}}{ \longleftrightarrow}}
\newcommand{\sinc}{\,\text{sinc}\,}
\newcommand{\rect}{\,\text{rect}\,}
	%\newcommand{\solution}[2]{\textbf{Solution:}{#1}}
\newcommand{\solution}{\noindent \textbf{Solution: }}
\newcommand{\cosec}{\,\text{cosec}\,}
\providecommand{\dec}[2]{\ensuremath{\overset{#1}{\underset{#2}{\gtrless}}}}
\newcommand{\myvec}[1]{\ensuremath{\begin{pmatrix}#1\end{pmatrix}}}
\newcommand{\mydet}[1]{\ensuremath{\begin{vmatrix}#1\end{vmatrix}}}
\newcommand*{\permcomb}[4][0mu]{{{}^{#3}\mkern#1#2_{#4}}}
\newcommand*{\perm}[1][-3mu]{\permcomb[#1]{P}}
\newcommand*{\comb}[1][-1mu]{\permcomb[#1]{C}}
%\numberwithin{equation}{section}
\numberwithin{equation}{section}
%\numberwithin{problem}{section}
%\numberwithin{definition}{section}
\makeatletter
\@addtoreset{figure}{problem}
\makeatother

\let\StandardTheFigure\thefigure
\let\vec\mathbf
%\renewcommand{\thefigure}{\theproblem.\arabic{figure}}
\renewcommand{\thefigure}{\theproblem}
%\setlist[enumerate,1]{before=\renewcommand\theequation{\theenumi.\arabic{equation}}
%\counterwithin{equation}{enumi}


%\renewcommand{\theequation}{\arabic{subsection}.\arabic{equation}}

\vspace{3cm}


%\usepackage{babel}
\begin{document}

\begin{tabular}{L{6cm} C{2cm} R{5cm} }

	\definecolor{circleorange}{rgb}{1,0.17,0.08}
\definecolor{darkorange}{rgb}{1,0.27,0.1}
\definecolor{orange2}{rgb}{1,0.5,0.15}
\definecolor{orange3}{rgb}{1,0.65,0.25}
\definecolor{yellow1}{rgb}{0.95,0.77,0.2}
%\begin{tikzpicture}[scale=0.2,every node/.style={transform shape}]
\begin{tikzpicture}[scale=0.1,every node/.style={transform shape}]
\draw [fill=circleorange,circleorange] (5,10) circle (1.15); 
\fill [darkorange] (5.06,8) -- (5.06,2) -- (7.3,1.2) -- (7.3,8.8) -- (5.06,8);
\fill [darkorange] (4.94,8) -- (4.94,2) -- (2.7,1.2) -- (2.7,8.8) -- (4.94,8);
\fill [orange2]    (7.4,8.4) -- (7.4,1.6) -- (8.2,1.2) -- (8.2,8.8) -- (7.4,8.4);
\fill [orange2]    (2.6,8.4) -- (2.6,1.6) -- (1.8,1.2) -- (1.8,8.8) -- (2.6,8.4);
\fill [orange3]    (8.3,8.4) -- (8.3,1.6) -- (9.0,1.2) -- (9.0,8.8) -- (8.3,8.4);
\fill [orange3]    (1.7,8.4) -- (1.7,1.6) -- (1.0,1.2) -- (1.0,8.8) -- (1.7,8.4);
\fill [yellow1]    (9.1,8.4) -- (9.1,1.6) -- (9.7,1.2) -- (9.7,8.8) -- (9.1,8.4);
\fill [yellow1]    (0.9,8.4) -- (0.9,1.6) -- (0.3,1.2) -- (0.3,8.8) -- (0.9,8.4);
\ifxetex
\node [scale=2.1] at (5,-0.1)  {   {\bf {\nakulafont  भारतीय प्रौद्योगिकी संस्थान हैदराबाद }} };
\node [scale=1.8] at (5,-1.2) {   {\bf { Indian Institute of Technology Hyderabad}} };
%\node [scale=1.8] at (5,-1.2) {   {\bf {\liberationsansfont Indian Institute of Technology Hyderabad}} };
\fi
\end{tikzpicture}
% \includegraphics[scale=0.05]{logo_iith} \newline

& Quiz 4 
	&
EE3900	
\end{tabular}


\vspace{-6mm}
\begin{center}
%\includegraphics[scale=0.95]{Yellow-Line}
\begin{tikzpicture}
\definecolor{yellow1}{rgb}{0.95,0.77,0.2}
\draw[line width=0.75mm, yellow1] (0,0) -- (\textwidth,0);
\end{tikzpicture}
\par\end{center}

 \section{Definitions}
\begin{enumerate}[label=\arabic*.,ref=\thesection.\theenumi]
\numberwithin{equation}{section}
%\item 
%\begin{align}
%u(n) =
%\begin{cases}
%1 & n \ge 0
%\\
%0 & n < 0
%\end{cases}
%\end{align}
\item The convolution  of 
	$ p_{X_1}(n)$ and $p_{X_2}(n)$ is defined as 
\begin{align}
	 p_{X_1}(n)*p_{X_2}(n) = 
	\sum_{k}^{}p_{X_1}(n-k)p_{X_2}(k)  
\end{align}
%The $Z$-transform of $p_X(n)$ is defined as 
%%\cite{proakis_dsp}
%\begin{align}
%P_X(z) = \sum_{n = -\infty}^{\infty}p_X(n)z^{-n}, \quad z \in \mathbb{C}
%\label{eq:dice_xz}
%\end{align}
%\item Let $X(t)$ be a random process.  The {\em autocorrelation function} is then defined as
%	\begin{align}
%		R_X\brak{\tau} = E\sbrak{X(t)X^*\brak{t+\tau}}
%	\end{align}
%\item The {\em Fourier} transform of $g(t)$ is defined as
%	\begin{align}
%		G(f) = \int_{-\infty}^{\infty}g(t)e^{-\j2\pi f t}\, dt
%	\end{align}
%\item The {\em power spectral density} of $X(t)$ is defined as
%	\begin{align}
%		R_X(\tau) &\system{F}	S_X(f) 
%	\end{align}
%i\item 
%	\begin{align}
%		g(t) &\system{F}	G(f) 
%		\\
%		\implies G(t) &\system{F} g(-f)
%	\end{align}
%\item 
%	\begin{align}
%		g(at) &\system{F}	\frac{1}{\abs{a}}G\brak{\frac{f}{\abs{a}}}
%	\end{align}
%\item 
%	\begin{align}
%		\rect(t) &= 
%		\begin{cases}
%			1 & t \in \brak{-\frac{1}{2}, \frac{1}{2}}
%			\\
%			0 & \text{otherwise}
%		\end{cases}
%		\\
%		\sinc(f) &= \frac{\sin\brak{\pi f}}{\pi f}
%	\end{align}
%\item If 
%	\begin{align}
%		Y(t) &= X(t)*h(t),
%		\\
%		S_Y(f) &= \abs{H(f)}^2 S_X(f)
%	\end{align}
%\item 
%	\begin{align}
%		\triangle\brak{\frac{t}{2}} &= 
%		\begin{cases}
%			1 -\abs{t}& t \in \brak{-1, 1}
%			\\
%			0 & \text{otherwise}
%		\end{cases}
%	\end{align}
%\item 
%	\begin{align}
%		g_1(t) * g_2(t) \system{F} G_1(f)G_2(f)
%	\end{align}
%\item The {\em Dirac delta} function is defined as 
%	\begin{align}
%		\delta(t) &= 0, t \ne 0 
%		\\
%		\int_{-\infty}^{\infty} \delta(t)\,dt &= 1
%	\end{align}
 \end{enumerate}
 \section{Problems}
\begin{enumerate}[label=\arabic*.,ref=\thesection.\theenumi]
\numberwithin{equation}{section}
\item 
Two dice, one blue and one grey, are thrown at the same time.   The event defined by the sum of the two numbers appearing on the top of the dice can have 11 possible outcomes 2, 3, 4, 5, 6, 6, 8, 9, 10, 11 and 12.  A student argues that each of these outcomes has a probability $\frac{1}{11}$.  Do you agree with this argument?  Justify your answer.
\item Let $X_i \in \cbrak{1,2,3,4,5,6}, i = 1,2,$ be the random variables representing the outcome for each die.  Assuming the dice to be fair, show that the probability mass function (pmf) is expressed as 
\begin{align}
\label{eq:dice_pmf_xi}
p_{X_i}(n) = \pr{X_i = n} = 
\begin{cases}
\frac{1}{6} & 1 \le n \le 6
\\
0 & otherwise
\end{cases}
\end{align}
\item Let 
\begin{align}
\label{eq:dice_xdef}
X &= X_1 + X_2,
%\\
%\implies X &\in \cbrak{1,2,\dots,12}
\end{align}
%
%The objective is to show that
%\begin{align}
%p_X(n) \ne \frac{1}{11}
%\label{eq:dice_wrong}
%\end{align}
\item Using the fact that 
%From \eqref{eq:dice_xdef},
\begin{align}
p_X(n) 
&= \sum_{k}^{}\pr{X_1  = n -k | X_2 = k}p_{X_2}(k)
%	&= \pr{X_1 + X_2 = n} 
%	\\
	%= \pr{X_1  = n -X_2}
%\\
\label{eq:dice_x_sum}
\end{align}
show that 
\begin{align}
p_X(n)  
	&= p_{X_1}(n)*p_{X_2}(n)
	%\sum_{k}^{}p_{X_1}(n-k)p_{X_2}(k) = 
\label{eq:dice_x_conv}
\end{align}
assuming  $X_1$ and $X_2$ to be independent,
%\begin{multline}
%\pr{X_1  = n -k | X_2 = k} 
%\\
%= \pr{X_1  = n -k} = p_{X_1}(n-k)
%\label{eq:dice_x1_indep}
%\end{multline}
%From \eqref{eq:dice_x_sum} and \eqref{eq:dice_x1_indep},
%where $*$ denotes the convolution operation. 
%%\cite{proakis_dsp}.  
%Substituting from \eqref{eq:dice_pmf_xi}
%in \eqref{eq:dice_x_conv},
%\begin{align}
%p_X(n) = \frac{1}{6}\sum_{k=1}^{6}p_{X_1}(n-k)= \frac{1}{6}\sum_{k=n-6}^{n-1}p_{X_1}(k)
%\label{eq:dice_x_conv_x1}
%\end{align}
%\begin{align}
%\because p_{X_1}(k) &= 0, \quad k \le 1, k \ge 6.
%\end{align}
%From \eqref{eq:dice_x_conv_x1},
%
\item Show that 
%\begin{align}
%p_X(n) &= 
%\begin{cases}
%0 & n < 1
%\\
%\frac{1}{6}\sum_{k=1}^{n-1}p_{X_1}(k) &  1 \le n-1 \le  6
%\\
%\frac{1}{6}\sum_{k=n-6}^{6}p_{X_1}(k) & 1 < n-6 \le 6
%\\
%0 & n > 12
%\end{cases}
%\label{eq:dice_x_conv_cond}
%\end{align}
%Substituting from \eqref{eq:dice_pmf_xi} in \eqref{eq:dice_x_conv_cond},
\begin{align}
p_X(n) &= 
\begin{cases}
0 & n < 1
\\
\frac{n-1}{36} &  2 \le n \le  7
\\
\frac{13-n}{36} & 7 < n \le 12
\\
0 & n > 12
\end{cases}
\label{eq:dice_x_conv_final}
\end{align}
%satisfying \eqref{eq:dice_wrong}.
%\item Show that 
%%
%%From \eqref{eq:dice_pmf_xi} and \eqref{eq:dice_xz}, 
%\begin{align}
%P_{X_1}(z) =P_{X_2}(z) 
%%	&= \frac{1}{6}\sum_{n = 1}^{6}z^{-n}
%%\\
%&=\frac{z^{-1}\brak{1-z^{-6}}}{6\brak{1-z^{-1}}}, \quad \abs{z} > 1
%\label{eq:dice_xiz}
%\end{align}
%upon summing up the geometric progression.  
%\item Show that 
%\begin{align}
%%\because p_X(n) &= p_{X_1}(n)*p_{X_2}(n),
%%\\
%P_X(z) &= P_{X_1}(z)P_{X_2}(z)
%\label{eq:dice_xzprod_def}
%\end{align}
%%The above property follows from Fourier analysis and is fundamental to signal processing. 
%%\cite{proakis_dsp}. 
%%From \eqref{eq:dice_xiz} and \eqref{eq:dice_xzprod_def},
%\item Show that 
%\begin{align}
%P_X(z) &= \cbrak{\frac{z^{-1}\brak{1-z^{-6}}}{6\brak{1-z^{-1}}}}^2
%\\
%&= \frac{1}{36}\frac{z^{-2}\brak{1-2z^{-6}+z^{-12}}}{\brak{1-z^{-1}}^2}
%\label{eq:dice_xzprod}
%\end{align}
%Using the fact that 
%%\cite{proakis_dsp}
%\begin{align}
%p_X(n-k) &\system{Z}P_X(z)z^{-k},
%\\
%nu(n)&\system{Z} \frac{z^{-1}}{\brak{1-z^{-1}}^2}
%\end{align}
%after some algebra, it can be shown that
%%{\tiny
%\begin{multline}
%\frac{1}{36}\lsbrak{\brak{n-1}u(n-1) - 2 \brak{n-7}u(n-7)}
%\\
%\rsbrak{ +\brak{n-13}u(n-13)}
%\\
%\system{Z}
%\frac{1}{36}\frac{z^{-2}\brak{1-2z^{-6}+z^{-12}}}{\brak{1-z^{-1}}^2}
%\label{eq:dice_xz_closed}
%\end{multline}
%%}
%
%where 
%
%From \eqref{eq:dice_xz}, \eqref{eq:dice_xzprod} and \eqref{eq:dice_xz_closed}
%\begin{multline}
%p_{X}(n) = \frac{1}{36}\lsbrak{\brak{n-1}u(n-1) 
%}
%\\
%\rsbrak{- 2 \brak{n-7}u(n-7)+\brak{n-13}u(n-13)}
%\end{multline}
%which is the same as \eqref{eq:dice_x_conv_final}.  Note that  \eqref{eq:dice_x_conv_final} can be obtained from \eqref{eq:dice_xz_closed} using contour integration as well.
%% \cite{proakis_dsp}.  

 \end{enumerate}

\end{document}
